\section{Comparación de los flujos completos de entrenamiento e inferencia}

Para ayudar a los lectores a establecer una comprensión integral, en esta sección se contrasta mediante tablas las diferencias entre los modelos difusivos y los modelos de flujo en cada etapa del entrenamiento y la inferencia.

\subsection{Comparación de la fase de entrenamiento}

\begin{table}[H]
\centering
\caption{Modelos difusivos vs.\ modelos de flujo: comparación de la fase de entrenamiento}
\begin{tabular}{@{} p{3.8cm} p{5cm} p{5cm} @{}}
\toprule
\textbf{Etapa} & \textbf{Modelo difusivo} & \textbf{Modelo de flujo (Rectified Flow)} \\
\midrule
Dominio temporal & $t \in \{0, \ldots, T\}$ o ${[}0, T{]}$ & $t \in {[}0, 1{]}$ \\
Proceso hacia adelante & $\mathbf{x}_t = \sqrt{\bar\alpha_t}\,\mathbf{x}_0 + \sqrt{1-\bar\alpha_t}\,\boldsymbol{\epsilon}$ & $\mathbf{x}_t = (1-t)\mathbf{x}_0 + t\mathbf{x}_1$ \\
Objetivo de predicción & Ruido $\boldsymbol{\epsilon}$ o gradiente logarítmico $\nabla\log p_t$ & Velocidad $v = \mathbf{x}_1 - \mathbf{x}_0$ \\
Función de pérdida de entrenamiento & $\|\boldsymbol{\epsilon}_\theta(\mathbf{x}_t,t) - \boldsymbol{\epsilon}\|^2$ & $\|v_\theta(\mathbf{x}_t,t) - (\mathbf{x}_1-\mathbf{x}_0)\|^2$ \\
Requerimientos de datos & Solo se necesita $\mathbf{x}_0 \sim p_{\text{data}}$ & Se requiere $\mathbf{x}_0 \sim p_0$ y $\mathbf{x}_1 \sim p_1$ \\
\bottomrule
\end{tabular}
\end{table}

\subsection{Comparación de la fase de inferencia}

\begin{table}[H]
\centering
\caption{Modelos difusivos vs.\ modelos de flujo: comparación de la fase de inferencia}
\begin{tabular}{@{} p{3.8cm} p{5cm} p{5cm} @{}}
\toprule
\textbf{Etapa} & \textbf{Modelo difusivo} & \textbf{Modelo de flujo} \\
\midrule
Punto inicial & $\mathbf{x}_T \sim \mathcal{N}(\mathbf{0},\mathbf{I})$ & $\mathbf{x}_0 \sim \mathcal{N}(\mathbf{0},\mathbf{I})$ \\
Método de resolución & Ecuación diferencial estocástica (SDE) o ODE con función de trayectoria (PF-ODE) & ODE \\
Actualización por paso & $\mathbf{x}_{t-1} = \mathbf{x}_t + f(\mathbf{x}_t,t)\Delta t$ & $\mathbf{x}_{t+\Delta t} = \mathbf{x}_t + v_\theta(\mathbf{x}_t,t)\Delta t$ \\
Número típico de pasos & 20--100 pasos & 1--10 pasos (después de Reflow) \\
Resolutores avanzados & DPM-Solver, entre otros & Euler / Midpoint son suficientes \\
\bottomrule
\end{tabular}
\end{table}

\subsection{Resumen de la sección}

Los modelos de flujo ofrecen un proceso más conciso tanto en la fase de entrenamiento como en la de inferencia. Es especialmente digno de notar que, después del paso Reflow, la fase de inferencia del Rectified Flow puede completarse con muy pocos pasos; esto representa su mayor ventaja para el despliegue práctico.
